% =============================================================
% บทที่ 3 - วิธีดำเนินงานโครงงาน
% หัวข้อ: การวิเคราะห์และทำนายผลการแข่งขันฟุตบอลลีกยุโรปด้วยกระบวนการวิทยาการข้อมูล
%        เพื่อการประเมินมูลค่าการเดิมพัน
% =============================================================

ในการดำเนินโครงงาน ``การวิเคราะห์และทำนายผลการแข่งขันฟุตบอลลีกยุโรปด้วยกระบวนการวิทยาการข้อมูลเพื่อการประเมินมูลค่าการเดิมพัน (European Football Match Outcome Prediction and Betting Market Analysis Using Data Science)'' ผู้จัดทำได้ออกแบบและพัฒนาท่อส่งข้อมูลเชิงโมดูลาร์ (Modular Data Pipeline) ตามกระบวนการวิทยาการข้อมูล เพื่อให้ระบบสามารถประมวลผลข้อมูลที่มีความซับซ้อน ป้องกันปัญหาข้อมูลอนาคตรั่วไหล (Data Leakage) และรองรับการประเมินผลกำไรทางการเงินร่วมกับการทำนายทางสถิติได้อย่างเป็นระบบ โดยแบ่งขั้นตอนการดำเนินงานออกเป็น 6 ส่วนหลัก ดังนี้

\begin{enumerate}[label=\arabic*)]
    \item การวิเคราะห์และกำหนดแนวทางการพัฒนาโครงงาน
    \item การเก็บรวบรวมข้อมูล (Data Collection)
    \item การทำความสะอาดข้อมูลและการตรวจสอบความถูกต้อง (Data Cleaning, Schema Validation \& Entity Resolution)
    \item การวิศวกรรมฟีเจอร์และการวิเคราะห์ตลาดเดิมพัน (Feature Engineering \& Betting Market Analysis)
    \item การพัฒนาแบบจำลอง การเปรียบเทียบกับตลาด และระบบทดสอบย้อนหลัง (Model Development, Market Test \& Backtesting Framework)
    \item การทดสอบและการประกันคุณภาพระบบ (Testing \& Quality Assurance Framework)
\end{enumerate}

% -----------------------------------------------------------
\section{การวิเคราะห์และกำหนดแนวทางการพัฒนาโครงงาน}
% -----------------------------------------------------------

การทำนายผลการแข่งขันฟุตบอลและการประเมินความคุ้มค่าทางการเดิมพันมีความซับซ้อนสูง เนื่องจากต้องคำนึงถึงทั้งปัจจัยเชิงสถิติของทีมและการตั้งราคาต่อรองของตลาด โครงงานนี้จึงมุ่งเน้นการใช้ประโยชน์จากข้อมูลที่ทราบก่อนเริ่มการแข่งขันเท่านั้น โดยมีแนวทางการดำเนินงานดังนี้

\begin{enumerate}[label=\arabic*)]
    \item \textbf{ศึกษาทฤษฎีและงานวิจัยที่เกี่ยวข้อง:} ศึกษาหลักการคำนวณระบบประเมินศักยภาพเชิงสัมพัทธ์ (Elo Rating System) \cite{hvattum2010using}, ค่าความคาดหวังในการทำประตู (Expected Goals: xG) \cite{understat2024}, แบบจำลองปัวซงสำหรับจำนวนประตู \cite{dixon1997modelling}, การคำนวณดัชนีความเหนื่อยล้าสะสม (Fatigue Index), การคำนวณความน่าจะเป็นของตลาดจากอัตราต่อรอง (Implied Probabilities) และกลยุทธ์การจัดสรรเงินทุน (Kelly Criterion \& Staking Strategies)
    \item \textbf{กำหนดขอบเขตและเป้าหมายการทำนาย:} กำหนดเป้าหมายการทำนายผลการแข่งขันออกเป็น 3 ผลลัพธ์ ได้แก่ เจ้าบ้านชนะ (Home Win: H), เสมอ (Draw: D) และทีมเยือนชนะ (Away Win: A) โดยทำนายเป็นความน่าจะเป็นของแต่ละผลลัพธ์ เพื่อนำไปค้นหาการเดิมพันที่มีมูลค่าคาดหมายเป็นบวก (Positive Expected Value: +EV)
    \item \textbf{กำหนดเครื่องมือและสภาพแวดล้อมการพัฒนา:} พัฒนาระบบด้วยภาษา Python 3.14 ขึ้นไป ร่วมกับเครื่องมือบริหารจัดการสภาพแวดล้อมเสมือน (\texttt{uv}), ไลบรารีวิเคราะห์ข้อมูลตาราง (Polars), เครื่องมือตรวจสอบสัญญาข้อมูล (Pandera และ Pydantic), ไลบรารีเรียนรู้ของเครื่อง (Scikit-learn), ระบบทดสอบซอฟต์แวร์อัตโนมัติ (Pytest) และระบบตรวจทานคุณภาพโค้ด (Ruff และ Pyright)
\end{enumerate}

% -----------------------------------------------------------
\section{การเก็บรวบรวมข้อมูล (Data Collection)}
% -----------------------------------------------------------

ผู้จัดทำได้พัฒนาระบบดึงข้อมูลอัตโนมัติ เพื่อรวบรวมข้อมูลจากแหล่งข้อมูลสาธารณะ โดยมีหลักการทำงานที่สำคัญ 2 ประการ คือ (1) สามารถเริ่มทำงานต่อได้ (Resumable) โดยข้ามไฟล์ที่มีอยู่แล้ว และ (2) ตรวจสอบความสมบูรณ์ก่อนบันทึก โดยเขียนข้อมูลลงไฟล์ชั่วคราวก่อน อ่านกลับมาตรวจกับโครงสร้างข้อมูลดิบ (Raw Schema) แล้วจึงเปลี่ยนชื่อเป็นไฟล์จริง เพื่อป้องกันไฟล์เสียหาย

\begin{enumerate}[label=\arabic*)]
    \item \textbf{ตารางการแข่งขันและสถิติ Understat:} ดึงผ่านไลบรารี soccerdata ครอบคลุม 5 ลีกใหญ่ยุโรป ตั้งแต่ฤดูกาล 2014--2015 ถึงฤดูกาล 2026--2027 เท่าที่แข่งไปแล้ว (21,882 นัด โดย 293 นัดเป็นของฤดูกาล 2026--2027 ถึงวันที่ 11 ตุลาคม 2026) ประกอบด้วยผลประตู ค่า xG/xA และสถิติรายบุคคลของผู้เล่นในแต่ละนัด \cite{understat2024}
    \item \textbf{Transfermarkt:} รวบรวมชุดข้อมูลปฏิทินการแข่งขันทุกรายการ (ลีก ถ้วยในประเทศ และรายการยุโรป) ของสโมสรใน 5 ลีกใหญ่ และตารางการลงสนามของผู้เล่น รวมถึงแผนการเล่น (Formation) ตัวจริง 11 คนแรก จากชุดข้อมูลสาธารณะ transfermarkt-datasets ซึ่งหยุดอัปเดตในเดือนกรกฎาคม 2026 จึงมีข้อมูลตั้งแต่ฤดูกาล 2014--2015 ถึง 2025--2026 เท่านั้น รวม 28,140 นัด แบ่งเป็นนัดในลีก 21,589 นัด ถ้วยในประเทศ 3,163 นัด รายการยุโรป 3,246 นัด และรายการอื่น 142 นัด \cite{transfermarkt2024}
    \item \textbf{Football-Data.co.uk:} ดึงข้อมูลสถิติผลการแข่งขันย้อนหลัง สถิติการยิง/เตะมุม และราคาอัตราต่อรอง (Betting Odds) จากเจ้ามือรับพนันหลายแห่ง (เช่น Bet365, Pinnacle) ย้อนหลังตั้งแต่ฤดูกาล 2014--2015 ถึงฤดูกาล 2026--2027 เท่าที่แข่งไปแล้ว สำหรับใช้ประเมินความน่าจะเป็นของตลาดและทดสอบย้อนหลัง (การเปรียบเทียบกับตลาดและชุดกลยุทธ์ใช้เฉพาะฤดูกาลที่จบแล้ว คือถึง 2025--2026) \cite{footballdata}
\end{enumerate}

% -----------------------------------------------------------
\section{การทำความสะอาดข้อมูลและการตรวจสอบความถูกต้อง}
% -----------------------------------------------------------

\subsection{ขั้นตอนการทำความสะอาด (Data Cleaning)}

กระบวนการทำความสะอาดแปลงไฟล์ข้อมูลดิบเป็นตารางสะอาดหลัก ได้แก่ ตารางนัดการแข่งขันลีก, ตารางสถิติผู้เล่นรายนัด, ตารางการแข่งขันทุกรายการของ Transfermarkt และตารางราคาต่อรองของ Football-Data.co.uk โดยมีขั้นตอนหลักดังนี้

\begin{enumerate}[label=\arabic*)]
    \item คัดเลือกเฉพาะนัดที่แข่งขันจบแล้ว และแปลงระบุฤดูกาลให้อยู่ในรูปแบบมาตรฐาน (เช่น 2014)
    \item แปลงชื่อสโมสรจากทุกแหล่งข้อมูลให้เป็นชื่อมาตรฐานเดียวกัน (ดูหัวข้อ Entity Resolution)
    \item แปลงรูปแบบแผนการเล่น (Formation) เช่น \texttt{4-2-3-1} เป็นตัวเลขจำนวนแนวรับ แนวกลาง และแนวรุก
    \item เชื่อมโยงนัดการแข่งขันระหว่าง Understat, Transfermarkt และ Football-Data.co.uk ตามวันที่และคู่แข่งขัน
\end{enumerate}

\subsection{การจับคู่และแปลงชื่อเอนทิตีให้เป็นมาตรฐาน (Entity Resolution)}

เนื่องจากชื่อสโมสรในแต่ละแหล่งข้อมูลมีความแตกต่างกัน ระบบจึงใช้พจนานุกรมเทียบเคียงชื่อ (Alias Table) แปลงเป็นชื่อมาตรฐาน (Canonical Names) ส่วนการจับคู่ระหว่างรหัสสโมสรของ Transfermarkt กับชื่อทีมของ Understat และ Football-Data.co.uk ใช้วิธีเรียนรู้จากตัวนัดการแข่งขัน โดยถือว่าเป็นนัดเดียวกันเมื่อแข่งขันในลีกเดียวกัน สกอร์ตรงกัน และวันที่แข่งขันห่างกันไม่เกิน 1 วัน วิธีนี้จับคู่สโมสรได้ 167 สโมสร ครอบคลุม 43,160 จาก 43,178 ครั้งที่สโมสรลงแข่ง (ร้อยละ 99.96) โดยสโมสรที่จับคู่ได้น้อยที่สุดยังครอบคลุมร้อยละ 97.4 ของนัดของสโมสรนั้น และไม่มีรหัสสโมสรใดถูกจับคู่กับสองชื่อ

\subsection{การตรวจสอบความถูกต้องของข้อตกลงข้อมูล (Data Schema Validation)}

ระบบประยุกต์ใช้แนวคิดสัญญาข้อมูล (Data Contracts) ด้วย Pandera โดยตรวจข้อมูลขาเข้าและขาออกในทุกขั้นตอน ร่วมกับโครงสร้างข้อมูลที่กำหนด แบ่งเป็น 3 ระดับ ดังนี้

\begin{itemize}
    \item \textbf{โครงสร้างข้อมูลดิบ (Raw Schema):} ตรวจไฟล์ข้อมูลดิบตอนดาวน์โหลดและอ่านเข้าสู่ระบบ
    \item \textbf{โครงสร้างข้อมูลระหว่างทาง (Step Schema):} ตรวจประเภทข้อมูลและขอบเขตค่าผลลัพธ์ในแต่ละขั้นตอนการประมวลผล
    \item \textbf{โครงสร้างข้อมูลสุดท้าย (Clean Schema):} ตรวจไฟล์ตารางสะอาดทุกไฟล์กับ Schema ที่กำหนดก่อนนำไปสร้างตัวแปร
\end{itemize}

% -----------------------------------------------------------
\section{การวิศวกรรมฟีเจอร์และการวิเคราะห์ตลาดเดิมพัน}
% -----------------------------------------------------------

สถาปัตยกรรมของระบบเป็นท่อส่งข้อมูลแบบโมดูลาร์ที่ประมวลผลตามลำดับ ดังตารางที่ \ref{tab:pipeline_architecture}

\begin{table}[htbp]
\centering
\caption{สถาปัตยกรรมโมดูลการประมวลผลข้อมูลในระบบ (Data Pipeline Modules)}
\label{tab:pipeline_architecture}
\small
\begin{tabular}{|c|l|>{\raggedright\arraybackslash}p{7.2cm}|}
\hline
\textbf{ลำดับ} & \textbf{คำสั่ง / โมดูล} & \textbf{หน้าที่การทำงานหลัก} \\
\hline
1 & \texttt{Data Collection} & ดึงข้อมูลดิบจาก Understat, Transfermarkt และ Football-Data.co.uk \\ \hline
2 & \texttt{Data Cleaning} & ทำความสะอาด แปลงชื่อสโมสร จับคู่นัดระหว่างแหล่งข้อมูล \\ \hline
3 & \texttt{Feature Engineering} & คำนวณ Elo, ความเหนื่อยล้า, ตารางแข่ง, ฟอร์มผู้เล่น และราคาต่อรองตลาด \\ \hline
4 & \texttt{Schema Check} & ตรวจสอบไฟล์ผลลัพธ์ทุกไฟล์กับ Pandera Data Contracts \\ \hline
5 & \texttt{Model \& Backtesting} & ฝึกแบบจำลอง ทำนายความน่าจะเป็น คำนวณ +EV และจำลองการวางเงินทุน \\ \hline
6 & \texttt{API Service} & ให้บริการตารางข้อมูล ผลการทำนาย และสัญญาณการเดิมพันผ่าน FastAPI \\
\hline
\end{tabular}
\end{table}

ตัวแปรทุกตัวใช้เฉพาะข้อมูลที่ทราบ ณ เวลาเริ่มการแข่งขัน โดยหน้าต่างค่าเฉลี่ยย้อนหลังไม่รวมนัดปัจจุบัน และคะแนน Elo ใช้ค่าก่อนแข่ง ผลลัพธ์รวมเป็นตารางตัวแปรหลักสำหรับใช้ฝึกแบบจำลอง

\subsection{การคำนวณคะแนนความแข็งแกร่งเชิงสัมพัทธ์ (Elo Rating Calculation)}
ระบบเล่นซ้ำผลการแข่งขันลีกทุกนัดตามลำดับวันที่ และบันทึกคะแนน Elo ของแต่ละทีมก่อนเริ่มนัด ($R_H$ และ $R_A$) โดยให้ผลต่างคะแนน $\Delta R = R_H - R_A$ ไม่รวมความได้เปรียบของเจ้าบ้าน ค่าคะแนนที่คาดหวังของเจ้าบ้าน ($E_H$) รวมความได้เปรียบของเจ้าบ้าน ($h$) และการปรับปรุงคะแนนหลังจบนัดคำนวณได้จากสมการที่ (\ref{eq:elo_expected}) และ (\ref{eq:elo_update})

\begin{equation}
    E_H = \frac{1}{1 + 10^{-(R_H + h - R_A)/400}}
    \label{eq:elo_expected}
\end{equation}

\begin{equation}
    R'_H = R_H + K \cdot G \cdot (S_H - E_H), \qquad R'_A = R_A - K \cdot G \cdot (S_H - E_H)
    \label{eq:elo_update}
\end{equation}

โดยที่ $S_H$ คือผลการแข่งขันจริงของเจ้าบ้าน ($1$ เมื่อชนะ, $0.5$ เมื่อเสมอ, $0$ เมื่อแพ้), $K = 15$ คือขนาดการปรับ, $h = 40$ คะแนน คือความได้เปรียบของเจ้าบ้าน และ $G$ คือตัวคูณตามผลต่างประตู $|gd|$ ซึ่งมีค่า $1$ เมื่อ $|gd| \le 1$, $1.5$ เมื่อ $|gd| = 2$ และ $(11 + |gd|)/8$ เมื่อ $|gd| \ge 3$ คะแนนเริ่มต้นเท่ากับ 1500 ทีมที่เลื่อนชั้นเริ่มที่ค่าเฉลี่ยคะแนนของทีมที่ตกชั้นในฤดูกาลก่อนหน้า ตัวอย่างขั้นตอนการคำนวณแสดงดังรายการที่ \ref{lst:elo_code}

\noindent
\begin{minipage}{\textwidth}
\begin{lstlisting}[language=Python, caption={ฟังก์ชันการคำนวณและปรับปรุงคะแนน Elo Rating}, label={lst:elo_code}]
def goal_multiplier(gd: int) -> float:
    gd = abs(gd)
    if gd <= 1:
        return 1.0
    if gd == 2:
        return 1.5
    return (11 + gd) / 8

def update_elo(r_home, r_away, home_goals, away_goals, k=15.0, home_adv=40.0):
    # expected score of the home side (home advantage included)
    expected = 1 / (1 + 10 ** (-(r_home + home_adv - r_away) / 400))
    # actual score: win 1, draw 0.5, loss 0
    if home_goals > away_goals:
        actual = 1.0
    elif home_goals == away_goals:
        actual = 0.5
    else:
        actual = 0.0
    delta = k * goal_multiplier(home_goals - away_goals) * (actual - expected)
    return r_home + delta, r_away - delta
\end{lstlisting}
\end{minipage}

\subsection{การคำนวณดัชนีความเหนื่อยล้าสะสมและความหนาแน่นตารางแข่ง}
ระบบคำนวณคะแนนความเหนื่อยล้า (Fatigue Score) ช่วง 0--100 จากภาระการเล่นของผู้เล่น 11 ตัวจริง (นาทีที่ลงเล่นใน 7/28 วัน และระยะพัก) แล้วรวมกับตัวแปรความหนาแน่นตารางแข่งขันระดับทีมจาก Transfermarkt ได้แก่ จำนวนวันพักนับจากนัดก่อนหน้า จำนวนนัดในรอบ 7 และ 14 วัน และตัวบอกการแข่งเกมเยือนในศึกยุโรป

\subsection{ฟอร์มรายบุคคลของผู้เล่นตัวจริง (Starter Form Features)}
สำหรับผู้เล่นตัวจริงแต่ละคน ระบบคำนวณค่า xG, xA และนาทีที่ลงเล่นใน 10 นัดก่อนหน้า คิดเป็นอัตราต่อ 90 นาที (\texttt{xg90}, \texttt{xa90}, \texttt{form90}) จากนั้นเรียงผู้เล่น 11 คนตามตำแหน่งไลน์อัปกระจายเป็นคอลัมน์ช่องตำแหน่ง (Slots) รวม 110 คอลัมน์

\subsection{การวิเคราะห์ราคาต่อรองอัตราจ่ายและมูลค่าการเดิมพัน (Odds \& Value Betting Analysis)}
นำราคาต่อรองจาก Football-Data.co.uk มาแปลงเป็นความน่าจะเป็นโดยนัยของตลาด (Implied Probabilities: $P_{market}$) พร้อมทั้งคำนวณค่าส่วนต่างเจ้ามือ (Overround / Vig) เพื่อปรับค่าความน่าจะเป็นให้รวมเป็น 100\% จากนั้นเปรียบเทียบกับความน่าจะเป็นที่ทำนายโดยแบบจำลอง ($P_{model}$) เพื่อหาค่าความคาดหมายเป็นบวก (Expected Value: $EV$):
\begin{equation}
    EV = (P_{model} \times \text{Odds}) - 1
\end{equation}
โดยนัดที่มีค่า $EV > 0$ จะถูกจัดเป็นโอกาสการเดิมพันที่มีมูลค่า (Value Bet Opportunity)

% -----------------------------------------------------------
\section{การพัฒนาแบบจำลอง การเปรียบเทียบกับตลาด และระบบทดสอบย้อนหลัง}
% -----------------------------------------------------------

\subsection{การคัดเลือกตัวแปร (Feature Selection)}
ตัวแปรที่แบบจำลองใช้ได้ถูกกำหนดเป็นรายชื่อกลุ่มที่อนุญาต (Allowlist) แต่ละกลุ่มประกอบด้วยค่าของทั้งสองทีมและผลต่าง ผลลัพธ์หลังจบเกม เช่น ประตูและ xG ของนัดนั้น จึงไม่สามารถรั่วเข้าสู่แบบจำลองได้ การคัดเลือกทำแบบเดินหน้า (Greedy Forward Selection) โดยเริ่มจากกลุ่มคะแนน Elo แล้วเพิ่มกลุ่มที่ลด Log Loss แบบ Walk-forward บนฤดูกาลตรวจสอบ 2021--2023 ได้มากที่สุดทีละกลุ่ม จนไม่มีกลุ่มใดลดได้เกินเกณฑ์ขั้นต่ำ ชุดทดสอบถูกใช้เพียงครั้งเดียวในตอนท้าย เพื่อเปรียบเทียบชุดที่คัดเลือกกับการใช้ Elo อย่างเดียวและการใช้ทุกกลุ่ม

\subsection{แบบจำลองทำนายผลการแข่งขัน}
แบบจำลองทุกแบบใช้ชุดตัวแปรที่คัดเลือกชุดเดียวกัน ปรับมาตรฐานตัวแปร และเติมค่าว่างด้วยค่ามัธยฐานของชุดฝึก ดังนี้
\begin{enumerate}[label=\arabic*)]
    \item \textbf{Multinomial Logistic Regression:} แบบจำลองหลัก ทำนายความน่าจะเป็น H/D/A โดยตรง
    \item \textbf{Logistic Regression + Sigmoid Calibration:} แบบจำลองเดียวกันที่ปรับเทียบความน่าจะเป็นด้วย Sigmoid แบบ 5-fold บนฤดูกาลฝึก
    \item \textbf{Gradient-boosted Trees:} ใช้ Scikit-learn \cite{scikit-learn} คลาส \texttt{HistGradientBoostingClassifier} เพื่อตรวจว่ามีความสัมพันธ์ไม่เชิงเส้นที่ Logistic Regression จับไม่ได้หรือไม่
    \item \textbf{Poisson Goals:} Poisson Regression หนึ่งตัวต่อหนึ่งฝั่งสำหรับจำนวนประตู แล้วรวมตารางสกอร์ที่เป็นอิสระต่อกันเป็นความน่าจะเป็น H/D/A ตามแนวคิดของ Dixon และ Coles \cite{dixon1997modelling}
\end{enumerate}
นอกจากนี้ยังสร้างเส้นโค้งการเรียนรู้ (Learning Curve) โดยฝึกแต่ละแบบจำลองด้วยนัดล่าสุดจำนวน 250 นัดจนถึงทั้งหมด แล้ววัดผลบนชุดทดสอบเดียวกัน เพื่อดูว่าการเพิ่มข้อมูลยังช่วยหรือไม่ และแบบจำลองจำข้อมูล (Overfitting) หรือไม่

\subsection{การแบ่งชุดข้อมูลตามลำดับเวลา (Temporal Split)}
ประเมินผลโดยแบ่งชุดข้อมูลตามลำดับเวลาจริง (ไม่สุ่ม) ฤดูกาลก่อนปี 2024 เป็นชุดฝึก ฤดูกาล 2021--2023 ใช้ตรวจสอบแบบ Walk-forward ระหว่างคัดเลือกตัวแปร (แต่ละฤดูกาลทำนายจากแบบจำลองที่ฝึกด้วยฤดูกาลก่อนหน้าเท่านั้น) และฤดูกาล 2024--2025 กับ 2025--2026 ที่แข่งจบแล้วเป็นชุดทดสอบที่ถูกใช้เพียงครั้งเดียว (ฤดูกาลที่ยังแข่งไม่จบไม่ถูกใช้วัดผล) เพื่อจำลองสถานการณ์ทำนายอนาคตตามจริงและป้องกันปัญหาข้อมูลรั่วไหล

\subsection{การเปรียบเทียบกับตลาดเดิมพัน (Market Test)}
เกณฑ์ที่แบบจำลองต้องเอาชนะคือความน่าจะเป็นจากราคาของ Pinnacle ซึ่งเป็นเจ้ามือที่ราคาแม่นที่สุดในบรรดาที่เผยแพร่ หลังหักส่วนต่างของเจ้ามือด้วยวิธี Power (เปรียบเทียบกับวิธี Proportional และ Shin) โดยใช้ราคา 2 ช่วงเวลา คือ ราคาเปิด (1--2 วันก่อนแข่ง ซึ่งผู้เดิมพันนำไปใช้ได้จริง) และราคาปิด (ณ เวลาเริ่มแข่ง) ผู้ทำนาย 5 แบบ ได้แก่ ราคาเปิด ราคาปิด ราคาเปิดที่ปรับเทียบใหม่ด้วย Logistic Regression ราคาเปิดร่วมกับฟอร์มทีม และฟอร์มทีมอย่างเดียว ถูกวัดผลบนนัดชุดเดียวกัน ด้วย Log Loss และ RPS แล้วเปรียบเทียบเป็นคู่รายนัด (Paired Difference) พร้อมช่วงความเชื่อมั่น 95\% แบบจำลองจะถือว่าชนะตลาดก็ต่อเมื่อช่วงความเชื่อมั่นอยู่ต่ำกว่าศูนย์ทั้งหมด

\subsection{กรอบการทดสอบย้อนหลังและชุดกลยุทธ์ (Backtesting Framework)}
การประเมินทางการเงินแบ่งเป็น 3 ส่วน ดังนี้
\begin{enumerate}[label=\arabic*)]
    \item \textbf{การทดสอบย้อนหลังของแบบจำลอง:} ทำนายแต่ละฤดูกาลตั้งแต่ 2017 ด้วย Logistic Regression แบบ Walk-forward แล้วเลือกผลที่มี $EV$ สูงสุดในแต่ละนัด วางเดิมพัน 1 หน่วยที่ราคาปิดของ Pinnacle เมื่อ $EV$ สูงกว่าเกณฑ์ $0\%$, $2\%$, $5\%$, $10\%$ และ $20\%$
    \item \textbf{ชุดกลยุทธ์ที่ลงทะเบียนล่วงหน้า (Pre-registered Strategy Suite):} กลยุทธ์ 18 กลยุทธ์ (42 รูปแบบ) ถูกเขียนในไฟล์ \texttt{strategies.toml} และบันทึกเข้า Git ก่อนรันกับข้อมูล ครอบคลุมกลยุทธ์พื้นฐาน (แทงทีมต่อ ทีมรอง ช่วงราคา และผลเดิมทุกนัด) กลยุทธ์ตามแบบจำลอง (Model EV และ Model Draw) และ Value Betting ที่ราคาดีที่สุดเทียบกับ Pinnacle \cite{kaunitz2017beating} ในตลาด 1X2, Asian Handicap และ Over/Under 2.5 แต่ละกลยุทธ์วัด ROI และกำไรคาดหมาย ณ ราคาปิด ($p_{close} \times \text{Odds} - 1$) พร้อมช่วงความเชื่อมั่น 95\% แยกรายลีก แล้วปรับแก้ $p$-value ของทั้ง 252 ช่องด้วยวิธี Benjamini--Hochberg \cite{benjamini1995controlling} และทดสอบการเลือกพารามิเตอร์แต่ละฤดูกาลจากฤดูกาลก่อนหน้าเท่านั้น
    \item \textbf{การบริหารเงินทุน (Staking):} เล่นซ้ำเดิมพันของกลยุทธ์ตัวแทนตามลำดับวันที่ ด้วยแผนวางเงิน 5 แบบ ได้แก่ Flat (1\% ของทุนเริ่มต้น), Kelly, Half Kelly, Quarter Kelly และ Martingale วัดเงินทุนสุดท้ายและ Maximum Drawdown แล้วสุ่มเดิมพันแบบ Bootstrap สร้างประวัติจำลอง 1,000 ชุด เพื่อวัดโอกาสเสียเงินทุนครึ่งหนึ่งและร้อยละ 90
\end{enumerate}
สุดท้ายประเมินความเป็นไปได้ในทางปฏิบัติของกลยุทธ์ที่มีความได้เปรียบ ได้แก่ จำนวนฤดูกาลที่ต้องใช้กว่าความได้เปรียบจะปรากฏเป็นกำไร การจำลองการปิดบัญชีของเจ้ามือ และผลตอบแทนบน Betfair Exchange หลังหักค่าคอมมิชชัน

\subsection{การทดสอบสมมติฐานทางสถิติ}
ทดสอบข้อสังเกตเชิงปริมาณด้วยระดับนัยสำคัญ 0.05 บนข้อมูลพรีเมียร์ลีกฤดูกาล 2000--2001 ถึง 2024--2025 ซึ่งยาวกว่าช่วงข้อมูลหลักเพื่อให้มีจำนวนนัดมากพอ (ตัดนัดใน 28 วันแรกของแต่ละฤดูกาล) ได้แก่ (1) ทดสอบความได้เปรียบของทีมเจ้าบ้านด้วย Paired t-test และ Wilcoxon Signed-rank Test แบบทางเดียว และ (2) ทดสอบความสัมพันธ์ระหว่างผลต่างคะแนน Elo กับผลต่างประตูด้วยสัมประสิทธิ์สหสัมพันธ์ของ Pearson และ Spearman

% -----------------------------------------------------------
\section{การทดสอบและการประกันคุณภาพระบบ (Testing \& Quality Assurance Framework)}
% -----------------------------------------------------------

\subsection{การทดสอบซอฟต์แวร์อัตโนมัติ (Automated Testing)}
จัดทำชุดทดสอบอัตโนมัติด้วย Pytest ครอบคลุมการคำนวณ Elo, ดัชนีความเหนื่อยล้า, ตัวแปรตารางแข่ง, การแปลงชื่อสโมสร, การตรวจสอบ Schema, การคำนวณ +EV, ระบบ Backtesting และบริการ API โดยใช้ Hypothesis สำหรับการทดสอบเชิงคุณสมบัติ (Property-based Testing)

\subsection{การตรวจสอบคุณภาพโค้ดและการเผยแพร่}
ใช้ Ruff ตรวจสอบรูปแบบโค้ด Pyright ตรวจสอบชนิดข้อมูล pre-commit ตรวจสอบก่อนบันทึกโค้ด และ GitHub Actions สำหรับรันการตรวจสอบและชุดทดสอบอัตโนมัติทุกครั้งที่มีการส่งโค้ดขึ้นระบบ พร้อมรองรับการใช้งานผ่าน Docker สภาพแวดล้อมเสมือน